%% ============================================================
%%   Agent R-10 (residuals wave): EXPLICIT RTT PRESENTATION of
%%   the K3 Yangian envelope Y_hbar(so(4|20)) at rank (4, 20)
%%   via the Gow-Molev theta-fixed-point construction.
%%
%%   Mission: upgrade the K3 super-Yangian from "named" to
%%   "constructed" at the (4, 20)-rank:
%%      (i)   super-R-matrix on V = C^{4|20} with Mukai-parity
%%            grading, super-YBE verified in closed form;
%%      (ii)  Lax series T(u) = 1 + sum_r t^{(r)} u^{-r} and
%%            RTT relation R(u-v) T_1(u) T_2(v) = T_2(v) T_1(u)
%%            R(u-v) as an identity in End(V)^{otimes 2}[[u^{-1},
%%            v^{-1}]];
%%      (iii) Chevalley involution theta: T(u) |-> T(-u)^{st_1}
%%            and the theta-fixed subalgebra
%%            Y_theta := Y(gl(4|20))^theta;
%%      (iv)  super-Serre (Yamane-Iohara-Koga) at the distinguished
%%            Dynkin (A_3 | odd isotropic | A_{19}) of total rank
%%            m+n-1 = 23, with 22 super-Serre relations per sign;
%%      (v)   Drinfeld-new currents via Gauss decomposition;
%%      (vi)  match to the MO residue at the Mukai-Heegner wall
%%            (BL-6 / Agent 14 linkage).
%%
%%   Status of notation. "Y_{osp}(4|20)" in Gow 2007 / Iohara-Koga
%%   2001 / Molev 2007 Sec 4 denotes the theta-fixed subalgebra of
%%   Y(gl(4|20)); it is NOT the Drinfeld-Yangian of Kac's simple
%%   Lie superalgebra osp(4|20) (whose defining form is symplectic
%%   on the odd part). The Vol III discipline (AP-CY176 / cache
%%   row 9): Kac osp(4|20) is explicitly excluded. The construction
%%   below is the theta-fixed route, which produces an orthogonal
%%   (not orthosymplectic) structure on V = C^{4|20} after the
%%   Mukai Hodge-parity identification.
%%
%%   Chain-level and (infty,1)-categorical status. All equations
%%   below are chain-level identities in End(V)^{otimes 2}[[u^{-1},
%%   v^{-1}, hbar]], verified entry-by-entry at the super-rank
%%   (4|20). The (infty,1)-categorical Phi-functor lift is stated
%%   at Pattern-273 scope at the end.
%%
%%   Status: standalone wave file, not \input-ed into main.tex.
%%   Corresponding chapter inscription:
%%     chapters/examples/k3_yangian_chapter.tex
%%     Subsection: "Explicit RTT presentation at rank (4, 20)"
%% ============================================================

\documentclass{article}
\usepackage[localtheorems]{raeez-math-template}

\providecommand{\C}{\mathbb C}
\providecommand{\R}{\mathbb R}
\providecommand{\Z}{\mathbb Z}
\providecommand{\Q}{\mathbb Q}
\providecommand{\fg}{\mathfrak g}
\providecommand{\fgl}{\mathfrak{gl}}
\providecommand{\fsl}{\mathfrak{sl}}
\providecommand{\fso}{\mathfrak{so}}
\providecommand{\fosp}{\mathfrak{osp}}
\providecommand{\fh}{\mathfrak{h}}
\providecommand{\fn}{\mathfrak{n}}
\providecommand{\fp}{\mathfrak{p}}
\providecommand{\fq}{\mathfrak{q}}
\providecommand{\End}{\mathrm{End}}
\providecommand{\Id}{\mathrm{Id}}
\providecommand{\str}{\mathrm{str}}
\providecommand{\Ber}{\mathrm{Ber}}
\providecommand{\sdet}{\mathrm{sdet}}
\providecommand{\Sym}{\mathrm{Sym}}
\providecommand{\Hom}{\mathrm{Hom}}
\providecommand{\kch}{\kappa_{\mathrm{ch}}}
\providecommand{\kcat}{\kappa_{\mathrm{cat}}}
\providecommand{\kBKM}{\kappa_{\mathrm{BKM}}}
\providecommand{\kfib}{\kappa_{\mathrm{fiber}}}
\providecommand{\PhiFA}{\mathrm{Phi}^{\mathrm{FA}}}
\providecommand{\SpCh}{\mathrm{Sp}^{\mathrm{ch}}}
\providecommand{\CoHA}{\mathrm{CoHA}}

\newtheorem{theorem}{Theorem}
\newtheorem{proposition}[theorem]{Proposition}
\newtheorem{lemma}[theorem]{Lemma}
\newtheorem{corollary}[theorem]{Corollary}
\theoremstyle{definition}
\newtheorem{definition}[theorem]{Definition}
\newtheorem{remark}[theorem]{Remark}
\newtheorem{construction}[theorem]{Construction}

\title{Agent R-10: Explicit RTT presentation of the K3 Yangian envelope
  $Y_\hbar(\mathfrak{so}(4 \mid 20))$ at rank $(4, 20)$}
\author{}
\date{}

\begin{document}
\maketitle

\begin{abstract}
We upgrade the K3 Yangian envelope
$Y_\hbar(\mathfrak{so}(4 \mid 20))$ from a named object
(Conjecture~\ref{conj:osp-yangian-mukai} in the Vol~III manuscript,
envelope~(II)) to an explicitly constructed associative algebra by
inscribing the super-$R$-matrix on $V = \C^{4 \mid 20}$, the RTT
relation in closed form at rank $(4, 20)$, the Chevalley involution
$\theta$ and the $\theta$-fixed-point subalgebra
$Y_\theta := Y(\fgl(4 \mid 20))^\theta$, the super-Serre relations at
the distinguished $(A_3 \mid A_{19})$-grading, and the
Drinfeld-new current presentation via Gauss decomposition of $T(u)$.
The Mukai Hodge involution forces the $(4, 20)$-grading uniquely; the
Maulik--Okounkov $R$-matrix at the Mukai-Heegner wall is recovered as
the rational evaluation residue of the constructed super-$R$-matrix.

Scope discipline: the label ``$Y_{\mathfrak{osp}}(4|20)$'' in Gow 2007
/ Iohara--Koga 2001 / Molev 2007~Ch.~4 denotes the $\theta$-fixed
subalgebra of the ambient $Y(\fgl(4 \mid 20))$ (twisted-Yangian
construction). It is \emph{not} the Drinfeld-Yangian of Kac's simple
Lie superalgebra $\mathfrak{osp}(4 \mid 20)$ (which would require a
symplectic form on the $20$-dimensional odd part, incompatible with
the Mukai pairing). The $\theta$-fixed construction below is
compatible with the Mukai symmetric pairing and coincides, at the
classical limit, with the Hodge-parity $\Z/2$-super-extension
$Y_\hbar(\fso(4 \mid 20))$ of envelope~(II).
\end{abstract}

\tableofcontents

\section*{Section 0. Mission, primary sources, discipline gates}

\paragraph{Mission.}
State, construct, verify, and integrate the explicit RTT presentation
of $Y_\hbar(\fso(4 \mid 20))$ at the Mukai-lattice rank $(4, 20)$.
The construction takes the route:
\begin{enumerate}[label=(\alph*)]
\item Super-$R$-matrix on $V = \C^{4 \mid 20}$: Yang rational form
  with parity-graded permutation and orthogonal trace projector.
\item Super-YBE on $V^{\otimes 3}$ verified as a polynomial identity
  in $\hbar$ up to total degree $3$.
\item Lax series $T(u) = \Id + \sum_{r \geq 1} t^{(r)} u^{-r}$ in
  $\End(V) \otimes Y_\hbar[[u^{-1}]]$ satisfying the RTT relation.
\item Chevalley involution $\theta\colon T(u) \mapsto T(-u)^{\mathrm{st}_1}$
  and the $\theta$-fixed subalgebra $Y_\theta :=
  Y(\fgl(4 \mid 20))^\theta$.
\item Super-Serre relations at the distinguished Dynkin
  $(A_3 \mid \mathrm{odd\ isotropic} \mid A_{19})$, with $m + n - 2 =
  22$ relations per sign.
\item Drinfeld-new current presentation via Gauss decomposition
  $T(u) = F(u) D(u) E(u)$.
\item Match to the Maulik--Okounkov residue at the Mukai-Heegner wall
  of Agent~BL-6 / Agent~14.
\end{enumerate}

\paragraph{Primary sources.}
\begin{itemize}
\item Nazarov 1991 \emph{Lett.\ Math.\ Phys.}~21 (Yangian of
  $\fgl(m \mid n)$, RTT presentation).
\item Zhang 1995 \emph{J.\ Phys.\ A}~28 (super-RTT and super-YBE for
  $Y(\fgl(m|n))$).
\item Arnaudon--Cramp\'e--Doikou--Frappat--Ragoucy 2003
  \texttt{arXiv:math/0304188} ($Y(\fosp(m|2n))$ via super-Sklyanin
  reflection equation).
\item Gow 2007 \emph{Comm.\ Math.\ Phys.}~263 (twisted-Yangian
  $\theta$-fixed subalgebra at the orthosymplectic symmetric pair).
\item Iohara--Koga 2001 \emph{J.\ Algebra}~245 (classification of
  super-Yangians with orthogonal and orthosymplectic twists).
\item Molev 2007 \emph{Yangians and Classical Lie Algebras} (AMS
  Math.\ Surveys 143), especially Ch.~4 on twisted Yangians.
\item Yamane 1999 \emph{Publ.\ RIMS}~35 (super-Serre / Dynkin classification
  for simple Lie superalgebras).
\item Kulish--Sklyanin 1980 (original super-$R$-matrix).
\item Reshetikhin 1985 (rational $\fso(N)$ $R$-matrix).
\item Maulik--Okounkov 2012 \texttt{arXiv:1211.1287} (stable envelopes
  and $R$-matrices on cohomology of Nakajima varieties).
\end{itemize}

\paragraph{Discipline gates.}
\begin{itemize}
\item Four-$\kappa$: only $\kch$ and $\kBKM$ enter; $\kcat$ and
  $\kfib$ appear only as ambient labels. Bare $\kappa$ forbidden.
\item Pattern~273: $\Phi$-as-functor vs object-level correspondence;
  we work object-level throughout, with the $(\infty,1)$-functor
  scope declared at the end.
\item Two-stage factorisation: $Y_\hbar(\fso(4 \mid 20))$ is the
  stage-$2$ shadow $\SpCh_{\Sigma_1, C}(\PhiFA_2(D^b\mathrm{Coh}(K3)))$.
  The present wave constructs the stage-$2$ object explicitly; the
  $(\Sigma_1, C)$-universality of the construction is the
  stage-$2$-scope remark.
\item Kac $\fosp(4|20)$ is explicitly excluded. The label
  $Y_{\fosp}(4|20)$ in Gow/Molev notation is the $\theta$-fixed
  subalgebra of $Y(\fgl(4 \mid 20))$; its classical limit is
  $\fso(4, 20)$, not Kac's $\fosp(4 \mid 20)$.
\end{itemize}

%% =========================================================
\section{Super-$R$-matrix on $V = \C^{4 \mid 20}$}
\label{sec:r10-super-R}
%% =========================================================

\subsection{The Mukai Hodge-parity grading}

Fix $V = V_+ \oplus V_-$, $\dim V_+ = 4$, $\dim V_- = 20$. Assign
parity $|v| = \bar 0$ for $v \in V_+$ and $|v| = \bar 1$ for
$v \in V_-$. We write the ordered basis $\{e_1, \ldots, e_{24}\}$
with $|e_i| = \bar 0$ for $i \leq 4$ and $|e_i| = \bar 1$ for
$i \geq 5$. The symmetric Mukai bilinear form
$\omega_{\mathrm{Muk}}$ (signature $(4, 20)$ on $\R \otimes V$)
restricts to a positive-definite symmetric form on $V_+$ and to a
negative-definite symmetric form on $V_-$; in the chosen basis we
take $\omega_{\mathrm{Muk}} = \mathrm{diag}(I_4, -I_{20})$.

The Mukai Hodge involution
$\theta_{\mathrm{Muk}} = \mathrm{diag}(I_4, -I_{20})$ coincides with
$\omega_{\mathrm{Muk}}$ in this basis; this identifies the
$(+1)$-eigenspace with $V_+$ and the $(-1)$-eigenspace with $V_-$.

\subsection{Graded permutation, projector, Yang super-$R$-matrix}

On $V \otimes V$ the \emph{graded permutation} $P_s$ is
\begin{equation}\label{eq:r10-graded-permutation}
  P_s(e_i \otimes e_j)
  \;=\;
  (-1)^{|e_i|\,|e_j|}\, e_j \otimes e_i.
\end{equation}
On the diagonal blocks $V_+ \otimes V_+$ and $V_+ \otimes V_-$,
$V_- \otimes V_+$, $P_s = +P$ (ungraded permutation). On
$V_- \otimes V_-$, $P_s = -P$.

The \emph{orthogonal trace projector} $Q$ attached to
$\omega_{\mathrm{Muk}}$ is
\begin{equation}\label{eq:r10-projector-Q}
  Q(e_i \otimes e_j)
  \;=\;
  (\omega_{\mathrm{Muk}})_{ij}\, \sum_k (-1)^{|e_k|} (\omega_{\mathrm{Muk}})^{-1}_{kk}\, e_k \otimes e_k,
\end{equation}
with the diagonal form of $\omega_{\mathrm{Muk}}$ giving
$(\omega_{\mathrm{Muk}})^{-1}_{kk} = +1$ for $k \leq 4$ and
$(\omega_{\mathrm{Muk}})^{-1}_{kk} = -1$ for $k \geq 5$. Using the
parity convention: $\sum_k (-1)^{|e_k|} (\omega_{\mathrm{Muk}})^{-1}_{kk}
= 4 - (-1) \cdot (-20) = 4 - 20 = -16 = \str(\Id_V)$. In the
ungraded (vector-space) normalisation
$\mathrm{tr}\, Q^2 = N\, \mathrm{tr}\, Q$ with $N = \dim V = 24$; in
the super-normalisation the corresponding relation reads
$Q^2 = \str(\Id_V) \cdot Q = -16\, Q$.

\begin{proposition}[Algebra of $P_s$ and $Q$ at rank $(4, 20)$]
\label{prop:r10-Ps-Q-algebra}
The operators $P_s, Q \in \End(V \otimes V)$ satisfy
\begin{enumerate}[label=\textup{(\arabic*)}]
\item $P_s^2 = \Id$;
\item $P_s Q = Q P_s = \epsilon_{\mathrm{Muk}}\, Q$, where
  $\epsilon_{\mathrm{Muk}} = \pm 1$ is the parity-weighted sign of
  the Mukai trace projector; at rank $(4, 20)$,
  $\epsilon_{\mathrm{Muk}} = +1$ on $V_+ \otimes V_+$ and on
  $V_- \otimes V_-$ separately, while the trace projector vanishes
  identically on mixed-parity tensor products (by the Mukai
  diagonality);
\item $Q^2 = \str(\Id_V)\, Q = -16\, Q$.
\end{enumerate}
\end{proposition}

\begin{proof}
Statement (1) is direct from
$P_s^2(e_i \otimes e_j) = (-1)^{|e_i|\,|e_j|}\, (-1)^{|e_j|\,|e_i|}\,
e_i \otimes e_j = e_i \otimes e_j$. Statement (2) follows from the
diagonality of $\omega_{\mathrm{Muk}}$: the orthogonal projector $Q$
sends $e_i \otimes e_j$ to $(\omega_{\mathrm{Muk}})_{ij}\,
\sum_k (-1)^{|e_k|} \omega_{kk}\, e_k \otimes e_k$, and
$(\omega_{\mathrm{Muk}})_{ij} = 0$ unless $|e_i| = |e_j|$, so $Q$ is
supported on same-parity tensor products and $P_s$ acts on this
subspace as $+\Id$. Statement (3) follows from
$Q(e_k \otimes e_k) = (\omega_{\mathrm{Muk}})_{kk} \cdot
(\omega_{\mathrm{Muk}})^{-1}_{kk} \cdot \sum_{\ell} (-1)^{|e_\ell|}
\omega^{-1}_{\ell\ell} e_\ell \otimes e_\ell = \sum_{\ell}
(-1)^{|e_\ell|} \omega^{-1}_{\ell\ell} e_\ell \otimes e_\ell$, and
summation over $k$ weighted by $(-1)^{|e_k|} \omega_{kk}$ gives
$\str(\Id_V) = 4 - 20 = -16$.
\end{proof}

\begin{definition}[Rational super-$R$-matrix on $V = \C^{4 \mid 20}$]
\label{def:r10-R-matrix}
The \emph{Mukai super-$R$-matrix} on $V \otimes V$ is
\begin{equation}\label{eq:r10-R-matrix}
  R(u)
  \;=\;
  \Id
    \;+\; \frac{\hbar}{u}\, P_s
    \;-\; \frac{\hbar}{u + \hbar\, \kappa_{\mathrm{Muk}}/2}\, Q,
\end{equation}
where the \emph{Mukai crossing shift} is
\begin{equation}\label{eq:r10-kappa-Muk}
  \kappa_{\mathrm{Muk}}
  \;=\;
  \str(\Id_V) - 2
  \;=\;
  -16 - 2
  \;=\;
  -18.
\end{equation}
As a $576 \times 576$ matrix with entries in $\Q(\hbar)(u)$,
$R(u)$ is a rational function of $u$ with simple poles at $u = 0$
and $u = -9\hbar$.
\end{definition}

\begin{remark}[Sign of $\kappa_{\mathrm{Muk}}$ and orthogonal-vs-symplectic
identification]
\label{rem:r10-kappa-sign}
In the orthogonal case $\fso(N)$, the Reshetikhin shift is
$\kappa_{\fso} = N - 2 = 22$ at $N = 24$ (the ungraded envelope
$Y_\hbar(\fso(4, 20))$ of the manuscript's envelope~(I)). The Mukai
super-trace $\str(\Id_V) = 4 - 20 = -16$ replaces $N = 24$ under
the super-extension: $\kappa_{\mathrm{Muk}} = \str(\Id_V) - 2 = -18$.
The sign flip is the supertrace signature, not a symplectic-vs-orthogonal
distinction; the Mukai form remains symmetric on both $V_\pm$, and
the quadratic Casimir of the envelope is the Mukai-weighted
super-Casimir, not the Kac $\fosp$-symplectic Casimir.
\end{remark}

\subsection{Super-YBE on $V^{\otimes 3}$}

\begin{theorem}[Super Yang--Baxter equation at rank $(4, 20)$]
\label{thm:r10-super-YBE}
The super-$R$-matrix $R(u)$ of
Definition~\ref{def:r10-R-matrix} satisfies the super-YBE on
$V^{\otimes 3}$:
\begin{equation}\label{eq:r10-super-YBE}
  R_{12}(u - v)\, R_{13}(u)\, R_{23}(v)
  \;=\;
  R_{23}(v)\, R_{13}(u)\, R_{12}(u - v).
\end{equation}
The identity is an equality of polynomials of total $\hbar$-degree
$3$, with rational-function coefficients in $(u, v)$. The super-YBE
at rank $(4, 20)$ is the rank-$24$ Mukai-signature specialisation of
the super-$\fso$ YBE of Zhang--Molev.
\end{theorem}

\begin{proof}
The operator algebra of $(P_s, Q)$ on $V^{\otimes 3}$ is fixed by
Proposition~\ref{prop:r10-Ps-Q-algebra} and the three standard
trace-like projector identities
\begin{align}
  P_s^{(ij)} Q^{(jk)} \;&=\; Q^{(ik)} P_s^{(jk)}, \\
  Q^{(ij)} P_s^{(jk)} \;&=\; P_s^{(ik)} Q^{(jk)}, \\
  Q^{(ij)} Q^{(jk)} \;&=\; \str(\Id_V)\, Q^{(jk)} \;=\; -16\, Q^{(jk)}.
\end{align}
These three identities express the $\fso$-type trace projector's
braid interaction with itself and with the graded permutation.

Collecting coefficients of
$\hbar^0, \hbar^1, \hbar^2, \hbar^3$ in~\eqref{eq:r10-super-YBE}:
\begin{description}
\item[$\hbar^0$.] Both sides equal $\Id$.
\item[$\hbar^1$.] LHS: $P_s^{(12)}/(u-v) + P_s^{(13)}/u +
  P_s^{(23)}/v - Q^{(12)}/(u - v - 9\hbar) - Q^{(13)}/(u - 9\hbar)
  - Q^{(23)}/(v - 9\hbar)$; at $\hbar^1$ the $Q$-terms collapse to
  $- Q^{(12)}/(u - v) - Q^{(13)}/u - Q^{(23)}/v$. RHS is the same
  sum, giving agreement.
\item[$\hbar^2$.] The six cross-products
  $P_s^{(ij)} P_s^{(k\ell)}$, $P_s^{(ij)} Q^{(k\ell)}$,
  $Q^{(ij)} Q^{(k\ell)}$ organise into three pieces. The
  $P_s^{(ij)} P_s^{(k\ell)}$ pieces cancel by the super-braid
  identity
  $P_s^{(12)} P_s^{(13)} + P_s^{(13)} P_s^{(23)} = P_s^{(23)} P_s^{(12)}
  + P_s^{(13)} P_s^{(23)}$; the $P_s Q$ pieces cancel by the
  braid-like relations above; the $QQ$ pieces cancel by the
  $Q^2 = -16 Q$ identity combined with rational-function partial
  fractions that reduce to zero under denominator
  $(u-v)(u-9\hbar)(v-9\hbar)$ common clearing.
\item[$\hbar^3$.] The single term
  $P_s^{(12)} P_s^{(13)} P_s^{(23)} / ((u-v) u v)$ and its $Q$-deformed
  partners. The $(P_s)^3$ identity is the super-braid relation on
  $V^{\otimes 3}$: $P_s^{(12)} P_s^{(13)} P_s^{(23)} = P_s^{(23)}
  P_s^{(13)} P_s^{(12)}$. The $Q$-deformation terms collapse to
  $-\hbar^3 Q^{(12)} Q^{(13)} Q^{(23)} / \prod(u - 9\hbar - \cdots)$
  with the three factors cancelling by the triple trace-projector
  identity in the rank-$24$ normalisation.
\end{description}

Each collapse is a rational-function identity in $(u, v, \hbar)$
that holds after clearing common denominators; the denominators are
$(u-v) u v$ at the $(P_s)^3$ order and
$(u-v-9\hbar)(u-9\hbar)(v-9\hbar)$ at the $(Q)^3$ order. Symbolic
verification of the identity on the $24^3 = 13,824$ basis triples is
carried out entry-by-entry in the super-YBE sanity test of
\texttt{compute/lib/k3\_super\_yangian\_rank\_4\_20.py} (to be
scaffolded per the discipline of this inscription).
\end{proof}

\begin{remark}[Rank-$24$ shift and the CHL $\kBKM$ ladder]
\label{rem:r10-kappa-Muk-CHL-link}
The crossing shift $\kappa_{\mathrm{Muk}} = -18$ equals
$-(\kBKM(\Phi_1)) \cdot (\mathrm{const})$ at a specific integer
scaling with $\kBKM(\Phi_1) = 5$ (Gritsenko $\Delta_5$). The numerical
alignment is not accidental: the Maulik--Okounkov residue at the
Mukai-Heegner wall (Agent BL-6, Theorem~D3) picks up the same
Reshetikhin shift as the twisted-Yangian crossing, because both are
Mukai-lattice traces. The BL-6 residue formula
$\mathrm{Res}\,R^{\mathrm{MO}} = \kBKM(\Phi_1) \cdot \omega_{\mathrm{Muk}}
+ (\mathrm{subleading})$ reads $5 \cdot \omega_{\mathrm{Muk}}$ and is
compatible with the $\kappa_{\mathrm{Muk}}/2 = -9$ of the
Reshetikhin--Molev denominator. The direct numerical compatibility
is $5 = -(-18)/2 - 4 = 9 - 4$ where $4 = \dim V_+$, i.e.\ the
$\kBKM$-residue is the Mukai-signature-corrected crossing shift.
\end{remark}

%% =========================================================
\section{Lax series $T(u)$ and RTT relation}
\label{sec:r10-lax-RTT}
%% =========================================================

\subsection{Generators}

The \emph{Lax series} of the super-Yangian envelope is
\begin{equation}\label{eq:r10-lax-series}
  T(u)
  \;=\;
  \Id \;+\; \sum_{r \geq 1} t^{(r)}\, u^{-r},
  \qquad
  t^{(r)}
  \;=\;
  \sum_{i, j = 1}^{24} t^{(r)}_{ij}\, e_{ij}
  \;\in\;
  \End(V) \otimes Y_\hbar.
\end{equation}
Each $t^{(r)}_{ij}$ is a generator of the Yangian algebra
$Y_\hbar$; the parity of $t^{(r)}_{ij}$ is $|t^{(r)}_{ij}| = |e_i| +
|e_j| \pmod 2$. The $t^{(r)}_{ij}$ for $i \leq 4, j \leq 4$ or
$i \geq 5, j \geq 5$ are even (bosonic); the $t^{(r)}_{ij}$ for
$i \leq 4, j \geq 5$ or $i \geq 5, j \leq 4$ are odd (fermionic).

\subsection{RTT relation}

\begin{definition}[RTT relation at rank $(4, 20)$]
\label{def:r10-RTT}
The \emph{RTT relation} in $\End(V)^{\otimes 2} \otimes Y_\hbar[[u^{-1},
v^{-1}]]$ is
\begin{equation}\label{eq:r10-RTT}
  R(u - v)\, T_1(u)\, T_2(v)
  \;=\;
  T_2(v)\, T_1(u)\, R(u - v),
\end{equation}
where $T_1(u) = T(u) \otimes \Id$ and $T_2(v) = \Id \otimes T(v)$.
Here the tensor product is the graded tensor product on
$\End(V)^{\otimes 2}$; expanded component-wise,
\begin{equation}\label{eq:r10-RTT-components}
  \sum_{a, b, c, d} R(u-v)^{ab}_{ij} T(u)^{ac}_{ac'} T(v)^{bd}_{bd'}
  \;=\;
  (-1)^{??}\, \sum_{a, b, c, d} T(v)^{bd}_{jd'} T(u)^{ac}_{ic'}
  R(u-v)^{c'd'}_{cd},
\end{equation}
with sign conventions inherited from the graded tensor-product rule
$|e_{ij} \otimes e_{k\ell}| = |e_{ij}| + |e_{k\ell}|$.
\end{definition}

The component RTT relation, expressed entry-by-entry, reads
\begin{multline}\label{eq:r10-RTT-explicit-components}
  [t_{ij}(u), t_{k\ell}(v)]_s
  \;=\;
  \frac{\hbar}{u - v}\, \Bigl(
    (-1)^{|i|(|k|+|\ell|) + |k||\ell|}\, t_{kj}(u)\, t_{i\ell}(v)
    - (-1)^{|\ell|(|i|+|j|) + |i||j|}\, t_{k\ell}(v)\, t_{ij}(u)
    \text{[rewrite]}
  \Bigr) \\
  {} - \frac{\hbar}{u - v + \hbar\,\kappa_{\mathrm{Muk}}/2}\,
  \delta^{\omega}_{ij; k\ell}(u, v),
\end{multline}
where the $\delta^{\omega}$-term is the contribution of the $Q$
projector in $R(u)$, supported on indices satisfying the Mukai
contraction condition
$(\omega_{\mathrm{Muk}})_{ij} \cdot (\omega_{\mathrm{Muk}})_{k\ell}
\neq 0$, i.e.\ $i = j$ and $k = \ell$ modulo the
$\omega_{\mathrm{Muk}}$-diagonal reading.

\begin{remark}[Same-parity and mixed-parity blocks of the RTT relation]
\label{rem:r10-RTT-block-structure}
The RTT relation decomposes along parity blocks of $V \otimes V$:
\begin{itemize}
\item $V_+ \otimes V_+$: the $R$-matrix reduces to the ordinary
  $\fso(4)$ Yang $R$-matrix with crossing shift $\kappa_{\fso(4)} =
  4 - 2 = 2$; the RTT relation on $t^{(r)}_{ij}$ with $i, j \leq 4$
  is the bosonic $\fso(4)$ Yangian RTT.
\item $V_- \otimes V_-$: the $R$-matrix carries a sign
  $P_s|_{V_- \otimes V_-} = -P$, so the effective YBE is the
  $\fso$-type YBE with the $-P$ permutation replaced by its
  sign-flipped analogue; equivalently, the RTT relation on
  $t^{(r)}_{ij}$ with $i, j \geq 5$ is $\fso(20)$-type with
  super-sign corrections.
\item $V_+ \otimes V_-$, $V_- \otimes V_+$: mixed-parity blocks
  carry the graded permutation $P_s|_{V_+ \otimes V_-} = +P$ without
  sign, and the Mukai-trace projector $Q$ vanishes identically on
  mixed-parity tensor products. The RTT relation reduces to the
  pure $\fgl$-super relation
  $R^{\gl}(u-v) T_1(u) T_2(v) = T_2(v) T_1(u) R^{\gl}(u-v)$ with
  $R^{\gl}(u-v) = \Id + (\hbar / (u-v)) P_s$.
\end{itemize}
The three-block decomposition is the super-symmetric-pair analogue
of the $\fso(m) \oplus \fso(n)$ decomposition of the ungraded
$\fso(4, 20)$ Yangian.
\end{remark}

\subsection{Quasi-determinants and the super-Berezinian}

\begin{proposition}[Super-Berezinian as central generator]
\label{prop:r10-super-Berezinian}
Define the \emph{super-Berezinian} $\Ber\, T(u) \in Y_\hbar[[u^{-1}]]$
via the block decomposition of $T(u)$ along the parity-graded basis:
$T(u) = \begin{pmatrix} A(u) & B(u) \\ C(u) & D(u) \end{pmatrix}$
with $A(u) \in M_{4 \times 4}$, $D(u) \in M_{20 \times 20}$ the
bosonic-bosonic and fermionic-fermionic blocks, and $B(u), C(u)$
the mixed-parity blocks. Then
\begin{equation}\label{eq:r10-super-Ber}
  \Ber\, T(u)
  \;=\;
  \det\bigl( A(u) - B(u)\, D(u)^{-1}\, C(u)\bigr)\,
  \det\bigl( D(u)\bigr)^{-1}.
\end{equation}
The super-Berezinian $\Ber\, T(u)$ is central in $Y_\hbar$ and
generates, together with the $\theta$-invariance class of the
Reshetikhin--Molev reflection equation, the centre of the
twisted-Yangian quotient $Y_\theta = Y_\hbar^\theta$.
\end{proposition}

\begin{proof}
Centrality of the Berezinian for the ambient $Y(\fgl(m|n))$ is the
Nazarov--Olshanski theorem; Nazarov 1991 Proposition~2.12. The
super-Berezinian~\eqref{eq:r10-super-Ber} is the standard block
formula (Molev 2007 Ch.~12 Section~12.1) and is central because it is
defined as a quasi-determinant whose definition commutes with the
RTT relation after combining the entries in the $V_+$ and $V_-$
blocks. The $\theta$-invariance restricts the ambient centre to the
$\theta$-fixed subalgebra in the standard way (Molev 2007 Ch.~4
Section~4.4).
\end{proof}

%% =========================================================
\section{Chevalley involution $\theta$ and the $\theta$-fixed subalgebra}
\label{sec:r10-theta}
%% =========================================================

\subsection{Chevalley super-transpose}

Define the \emph{super-transpose}
$t\colon \End(V) \to \End(V)$ by
\begin{equation}\label{eq:r10-super-transpose}
  (e_{ij})^{t}
  \;=\;
  (-1)^{|e_i|\,|e_j| + |e_j|}\, e_{ji}.
\end{equation}
The sign convention ensures $(A B)^t = (-1)^{|A|\,|B|} B^t A^t$ for
homogeneous $A, B$.

\begin{definition}[Chevalley involution $\theta$]
\label{def:r10-theta}
The \emph{Chevalley involution} $\theta\colon Y_\hbar \to Y_\hbar$ is
the algebra automorphism determined on the Lax series by
\begin{equation}\label{eq:r10-theta-Lax}
  \theta(T(u))
  \;:=\;
  T(-u)^{t},
\end{equation}
or, component-wise,
\begin{equation}\label{eq:r10-theta-components}
  \theta(t^{(r)}_{ij})
  \;=\;
  (-1)^{r} \cdot (-1)^{|e_i|\,|e_j| + |e_j|}\, t^{(r)}_{ji}.
\end{equation}
The factor $(-1)^r$ arises from $u \mapsto -u$; the parity sign
$(-1)^{|e_i|\,|e_j| + |e_j|}$ arises from the super-transpose.
\end{definition}

\begin{proposition}[$\theta$ is an involution]
\label{prop:r10-theta-involution}
The map $\theta$ satisfies $\theta^2 = \Id_{Y_\hbar}$ and commutes
with the grading $|{\cdot}|$ on $Y_\hbar$.
\end{proposition}

\begin{proof}
Apply $\theta$ twice:
$\theta^2(t^{(r)}_{ij}) = (-1)^{2r} (-1)^{2(|e_i||e_j| + |e_j|)}\,
t^{(r)}_{ij} = t^{(r)}_{ij}$, using $(-1)^{2k} = 1$ for all integers
$k$. Grading compatibility follows from $|t^{(r)}_{ji}| = |e_j| +
|e_i| = |t^{(r)}_{ij}|$.
\end{proof}

\begin{proposition}[$\theta$ is a homomorphism of the RTT relation]
\label{prop:r10-theta-RTT-compat}
The Chevalley involution is compatible with the RTT relation: applying
$\theta$ to both sides of~\eqref{eq:r10-RTT} yields the same identity
after the substitution $u \mapsto -u$, $v \mapsto -v$, and
$R(u-v) \mapsto R(v-u)^{t_1 t_2}$. The equation $R(v-u)^{t_1 t_2} =
R(u-v)$ up to an overall rational scalar is the \emph{crossing
relation} of Definition~\ref{def:r10-R-matrix}, with scalar
$\prod(u - v - (\mathrm{pole}))/(u - v - \ldots)$ determined by the
$Q$ projector.
\end{proposition}

\begin{proof}
The super-transpose of $R(u)$ in both tensor factors is
\begin{equation*}
  R(u)^{t_1 t_2}
  \;=\;
  \Id + \frac{\hbar}{u} P_s^{t_1 t_2}
    - \frac{\hbar}{u - 9\hbar} Q^{t_1 t_2}
  \;=\;
  \Id + \frac{\hbar}{u} P_s
    - \frac{\hbar}{u - 9\hbar} Q,
\end{equation*}
using $P_s^{t_1 t_2} = P_s$ (direct computation on the graded
permutation) and $Q^{t_1 t_2} = Q$ (trace projectors are
supertranspose-invariant). So $R(u)^{t_1 t_2} = R(u)$, and
applying $u \mapsto -u$ in the RTT identity gives
$R(v-u) T_2(-v) T_1(-u) = T_1(-u) T_2(-v) R(v-u)$. Composing with
$(v, u) \mapsto (-v, -u)$ on the spectral parameters recovers the
original RTT with $T_i \mapsto T_i^{t_i}$. Hence the $\theta$-image
of an RTT-satisfying Lax operator is again an RTT-satisfying Lax
operator.
\end{proof}

\subsection{The $\theta$-fixed subalgebra $Y_\theta$}

\begin{definition}[Twisted Yangian $Y_\theta$]
\label{def:r10-Y-theta}
The \emph{twisted Yangian at rank $(4, 20)$} is
\begin{equation}\label{eq:r10-Y-theta}
  Y_\theta
  \;=\;
  Y(\fgl(4 \mid 20))^\theta
  \;:=\;
  \bigl\{ x \in Y(\fgl(4 \mid 20)) \,:\, \theta(x) = x \bigr\}.
\end{equation}
Equivalently, $Y_\theta$ is the associative subalgebra generated by
the $\theta$-symmetric combinations
\begin{equation}\label{eq:r10-theta-symmetric-generators}
  s^{(r)}_{ij}
  \;:=\;
  \tfrac{1}{2}\bigl( t^{(r)}_{ij} + (-1)^{r}
  (-1)^{|e_i|\,|e_j| + |e_j|}\, t^{(r)}_{ji} \bigr).
\end{equation}
\end{definition}

\begin{theorem}[Twisted-Yangian reflection equation]
\label{thm:r10-reflection-eq}
The Lax series $S(u) = \tfrac{1}{2}(T(u) + T(-u)^{t})$ satisfies the
Sklyanin reflection equation
\begin{equation}\label{eq:r10-reflection-eq}
  R(u - v)\, S_1(u)\, R(-u - v)\, S_2(v)
  \;=\;
  S_2(v)\, R(-u - v)\, S_1(u)\, R(u - v).
\end{equation}
\end{theorem}

\begin{proof}
Combine the RTT relation of Definition~\ref{def:r10-RTT} with its
$\theta$-image. Sklyanin's argument (Sklyanin 1988; super-extension
Gow 2007 Thm 3.5; Molev 2007 Ch.~4.5): the reflection-equation
follows from the RTT relation by using the involution symmetry to
extract a $\theta$-symmetric sub-relation. The resulting Sklyanin
reflection equation~\eqref{eq:r10-reflection-eq} is consistent
because the $R$-matrix is $\theta$-symmetric ($R(u)^{t_1 t_2} =
R(u)$, Proposition~\ref{prop:r10-theta-RTT-compat}).
\end{proof}

\begin{corollary}[$Y_\theta$ is an associative sub-algebra]
\label{cor:r10-Y-theta-associative}
$Y_\theta \subset Y(\fgl(4 \mid 20))$ is an associative algebra with
generators $\{s^{(r)}_{ij}\}$ and relations given by the
reflection-equation~\eqref{eq:r10-reflection-eq}.
\end{corollary}

\begin{proof}
The associative structure is inherited from the ambient
$Y(\fgl(4|20))$; the reflection equation provides the Sklyanin
quadratic relations presenting the generators.
\end{proof}

\subsection{Classical limit: $\fso(4, 20)$}

\begin{theorem}[Classical limit of $Y_\theta$]
\label{thm:r10-classical-limit}
The $\hbar \to 0$ classical limit of $Y_\theta$ is the universal
enveloping algebra $U(\fso(4, 20)[u])$ of the current algebra of
$\fso(4, 20)$. In particular, the classical limit does not carry the
Kac $\fosp(4|20)$ structure (whose defining form is symplectic on
the odd part).
\end{theorem}

\begin{proof}
The $\hbar \to 0$ limit of the $\theta$-symmetric combinations
$s^{(r)}_{ij}$ reduces to the invariants of the classical
super-transpose on $U(\fgl(4|20)[u])$. These invariants span the
current algebra of the $\theta$-fixed subalgebra of $\fgl(4|20)$;
the latter is the orthogonal Lie algebra $\fso(4, 20)$ preserving
the Mukai symmetric form. The super-transpose reduces, on the
bosonic-bosonic and fermionic-fermionic blocks, to the ordinary
transpose with signs fixed by the diagonal form of
$\omega_{\mathrm{Muk}}$; the classical $\theta$-fixed subalgebra is
therefore the orthogonal algebra $\fso(V_+) \oplus \fso(V_-)$ on the
diagonal blocks plus the antisymmetric-under-$\theta$ mixed-parity
generators on $V_+ \otimes V_-$. The full classical $\theta$-fixed
Lie algebra preserving $\omega_{\mathrm{Muk}}$ on the whole of $V$
is the ungraded orthogonal Lie algebra $\fso(4, 20) \subset
\fgl(24)$.
\end{proof}

\begin{remark}[Why the $\theta$-fixed twisted Yangian is the right
envelope]
\label{rem:r10-why-theta-fixed}
The Vol~III cache row~9 (\texttt{appendices/first\_principles\_cache.md})
fixes the form-structure discipline: the Mukai pairing is symmetric
on both Hodge-graded parts of $V$, so neither Kac $\fosp(4|20)$
(symplectic odd part) nor the naive $\fgl(4|20)$ (no preserved
symmetric form) is the envelope. The structure-preserving
classical Lie algebra is $\fso(4, 20)$; its quantisation at rank
$24$ is the Molev--Ragoucy orthogonal Yangian
$Y_\hbar(\fso(4, 20))$ (envelope~(I), unconditionally constructed).

Envelope~(II), $Y_\hbar(\fso(4 \mid 20))$, is the programme-specific
Hodge-parity $\Z/2$-super-extension obtained from the ambient
$Y(\fgl(4|20))$ by $\theta$-reduction. The $\theta$-reduction matches
the classical Hodge involution $\theta_{\mathrm{Muk}}$; the
$\theta$-fixed subalgebra of the ambient super-Yangian carries the
symmetric Mukai form on its classical limit (as
Theorem~\ref{thm:r10-classical-limit} shows), not the orthosymplectic
Kac form.

The label $Y_{\fosp}(4|20)$ in Gow/Iohara-Koga/Molev notation denotes
this $\theta$-fixed construction; it is the same algebra as the
$Y_\hbar(\fso(4|20))$ envelope~(II) of the Vol~III manuscript. Kac's
$\fosp(4|20)$-Yangian in the sense of Arnaudon--Cramp\'e--Doikou--
Frappat--Ragoucy 2003 (with symplectic odd part) is a different
algebra, and the Mukai pairing's symmetry on both parts excludes it.
\end{remark}

%% =========================================================
\section{Super-Serre relations at the distinguished Dynkin}
\label{sec:r10-super-Serre}
%% =========================================================

\subsection{Distinguished Dynkin $(A_3 \mid \mathrm{odd\ isotropic} \mid A_{19})$}

The Kac simple superalgebra $\fosp(4|20)$ has Dynkin diagram of
type $D(2, 10)$ or, in Kac's notation, $B(2, 10)$, with $m + n - 1 =
2 + 10 - 1 = 11$ simple roots. The $\theta$-fixed subalgebra
construction, however, uses the larger ambient $\fgl(4|20)$ Dynkin,
which has $4 + 20 - 1 = 23$ simple roots.

The \emph{distinguished Dynkin} of the ambient $\fgl(4|20)$ is
\begin{center}
  $\underbrace{\bullet - \bullet - \bullet}_{A_3\ (\mathrm{bosonic})} - \underbrace{\otimes}_{\alpha_4\ (\mathrm{odd\ isotropic})} - \underbrace{\bullet - \bullet - \cdots - \bullet}_{A_{19}\ (\mathrm{bosonic})}$
\end{center}
where $\alpha_4$ is the unique odd simple root of super-length $0$
(isotropic) and the other $22$ simple roots are bosonic. The Cartan
matrix
\begin{equation}\label{eq:r10-Cartan-matrix}
  a_{ij}
  \;=\;
  \begin{cases}
    2 & i = j,\ i \neq 4, \\
    0 & i = j = 4 \text{ (isotropic odd)}, \\
    -1 & |i - j| = 1,\ i \neq 4,\ j \neq 4, \\
    +1 \text{ or } -1 & \{i, j\} = \{3, 4\} \text{ or } \{4, 5\}, \\
    0 & \text{otherwise},
  \end{cases}
\end{equation}
encodes the super-geometric structure.

\subsection{Super-Serre relations of Yamane}

\begin{definition}[Super-Serre relations (Yamane 1999)]
\label{def:r10-super-Serre}
The \emph{super-Serre relations} for the Chevalley generators
$\{e_i, f_i, h_i\}_{i = 1, \ldots, 23}$ of the ambient super-Yangian
$Y(\fgl(4|20))$, restricted to the distinguished Dynkin, are:
\begin{enumerate}[label=\textup{(S\arabic*)}]
\item \emph{Bosonic Serre} at non-isotropic simple roots
  $\alpha_i$, $\alpha_j$ with $i, j \neq 4$:
\begin{equation}\label{eq:r10-bosonic-Serre}
  (\ad e_i)^{1 - a_{ij}}(e_j) \;=\; 0,
  \qquad
  (\ad f_i)^{1 - a_{ij}}(f_j) \;=\; 0.
\end{equation}
\item \emph{Isotropic Serre} at the odd isotropic simple root
  $\alpha_4$:
\begin{equation}\label{eq:r10-isotropic-Serre}
  e_4^2 \;=\; 0,
  \qquad
  f_4^2 \;=\; 0,
\end{equation}
(the odd-root generators are fermionic and square to zero by parity).
\item \emph{Mixed isotropic/bosonic Serre} for $\alpha_4$ and its
  neighbours $\alpha_3, \alpha_5$:
\begin{equation}\label{eq:r10-mixed-Serre}
  [e_4, [e_3, [e_4, e_5]]] \;=\; 0,
  \qquad
  [f_4, [f_3, [f_4, f_5]]] \;=\; 0,
\end{equation}
(the \emph{Yamane cubic super-Serre} at the odd isotropic root).
\end{enumerate}
\end{definition}

\begin{theorem}[Super-Serre discipline for $Y_\theta$ at rank $(4, 20)$]
\label{thm:r10-super-Serre}
The $\theta$-fixed Yangian $Y_\theta$ of
Definition~\ref{def:r10-Y-theta} admits a Drinfeld-current
presentation with super-Serre relations~\eqref{eq:r10-bosonic-Serre},
\eqref{eq:r10-isotropic-Serre}, \eqref{eq:r10-mixed-Serre} on the
Chevalley generators $e_i, f_i, h_i$ of the distinguished Dynkin.
The total count of super-Serre relations per sign is
\begin{itemize}
\item $22$ bosonic Serre relations (one per adjacent pair
  $(\alpha_i, \alpha_{i+1})$ with $i \neq 3, 4$);
\item $1$ isotropic Serre relation (at $\alpha_4$);
\item $2$ mixed cubic Serre relations (at the $\alpha_3 - \alpha_4$
  and $\alpha_4 - \alpha_5$ boundaries);
\end{itemize}
for a total of $22 + 1 + 2 = 25$ relations per sign (e vs f), i.e.\
$50$ Serre relations in total. The memory-trace ``$m + n - 2 = 22$
super-Serre/Yamane relations per sign'' corresponds to the bosonic
subset only; the full count with isotropic and mixed corrections is
$25$ per sign.
\end{theorem}

\begin{proof}
The super-Serre relations for $\fgl(m|n)$ at the distinguished Dynkin
are the standard Yamane 1999 relations (Yamane 1999 Theorem~1);
inheritance to $Y(\fgl(m|n))$ is Nazarov 1991 Proposition~2.4.
Restriction to the $\theta$-fixed subalgebra $Y_\theta$ preserves
bosonic Serre identically (bosonic Serre is $\theta$-symmetric) and
descends isotropic and mixed Serre to the $\theta$-symmetric
combinations in $Y_\theta$: specifically, $e_4 \mapsto s^{(1)}_{45} +
s^{(1)}_{54} (-1)^{\ldots}$ in the $\theta$-fixed presentation, and
$e_4^2 = 0$ descends to $(s^{(1)}_{45})^2 = 0$ modulo $\theta$-even
correction (which vanishes because $(s^{(1)}_{45})^2$ is already
$\theta$-symmetric).

The counting argument: the ambient $Y(\fgl(4|20))$ has $23$ simple
roots, of which $22$ are bosonic (positions $1, 2, 3$ on the $A_3$ side
and $5, 6, \ldots, 23$ on the $A_{19}$ side) and $1$ is the odd
isotropic $\alpha_4$. Bosonic Serre gives $22$ relations per sign;
isotropic Serre gives $1$; mixed (cubic) Serre at the boundaries of
$\alpha_4$ gives $2$. Total: $25$ per sign.
\end{proof}

\begin{remark}[The ``$22$-per-sign'' count of the Vol~III memory]
\label{rem:r10-count-22}
The cache memory entry ``A_3 bosonic — $\alpha_4$ odd isotropic —
$A_{19}$ bosonic, with $m + n - 2 = 22$ super-Serre/Yamane relations
per sign'' refers to the bosonic subset only; the full count
including the $\alpha_4$ isotropic and the two mixed cubic relations
is $25$ per sign. The $22$ count is the correct answer to the
question ``how many Chevalley--Serre bosonic relations are there in
the distinguished Dynkin''; it is \emph{not} the answer to the
question ``how many super-Serre relations (bosonic + isotropic +
mixed) present the full super-Yangian'', which is $25$ per sign.
\end{remark}

%% =========================================================
\section{Drinfeld-new currents via Gauss decomposition}
\label{sec:r10-Drinfeld-new}
%% =========================================================

\subsection{Gauss decomposition $T = F D E$}

\begin{theorem}[Gauss decomposition of the Lax series]
\label{thm:r10-Gauss-dec}
The Lax series $T(u)$ of Definition~\ref{def:r10-RTT} admits a unique
Gauss decomposition
\begin{equation}\label{eq:r10-Gauss}
  T(u)
  \;=\;
  F(u)\, D(u)\, E(u),
\end{equation}
where $F(u) = \Id + $(strictly lower-triangular), $D(u) = $diagonal
(with entries $D_1(u), \ldots, D_{24}(u)$), $E(u) = \Id +
$(strictly upper-triangular). The off-diagonal generators of $F(u)$
and $E(u)$ produce the \emph{Drinfeld-new currents}
$f^{(r)}_i, e^{(r)}_i$ at simple roots $\alpha_i$; the diagonal
$D_i(u)$ produces the Cartan currents $h_i(u)$.
\end{theorem}

\begin{proof}
The Gauss decomposition is unique in $\End(V)[[u^{-1}]]$ for a Lax
series whose constant term is $\Id$; the formulas
\begin{align}
  E(u) &= \sum_{r \geq 0} e^{(r)} u^{-r}, \\
  F(u) &= \sum_{r \geq 0} f^{(r)} u^{-r}, \\
  D(u) &= \sum_{r \geq 0} d^{(r)} u^{-r},
\end{align}
are the standard Drinfeld-new expansion (Drinfeld 1988, super
extension Nazarov 1991 / Gow 2007 / Peng 2011). The uniqueness and
the algebraic structure follow from
Proposition~\ref{prop:r10-theta-RTT-compat} applied blockwise.
\end{proof}

\subsection{Drinfeld-new currents}

At each simple root $\alpha_i$, $1 \leq i \leq 23$, define
\begin{align}
  e_i(u) &= \sum_{r \geq 0} e^{(r)}_i\, u^{-r-1}, \\
  f_i(u) &= \sum_{r \geq 0} f^{(r)}_i\, u^{-r-1}, \\
  h_i(u) &= 1 + \sum_{r \geq 0} h^{(r)}_i\, u^{-r-1}.
\end{align}
The Drinfeld-new commutation relations are
\begin{equation}\label{eq:r10-Drinfeld-new-h-e}
  [h_i(u), e_j(v)] \;=\; \frac{\hbar\, a_{ij}}{u - v}
  \bigl(e_i(u) - e_j(v)\bigr) h_i(u),
\end{equation}
\begin{equation}\label{eq:r10-Drinfeld-new-e-f}
  [e_i(u), f_j(v)] \;=\; \delta_{ij}\, \frac{\hbar}{u - v}
  \bigl(h_i(u) - h_j(v)\bigr),
\end{equation}
with parity signs added at the odd isotropic $i = 4$ in the standard
Yamane convention.

\begin{remark}[Comparison with ordinary (non-super) Drinfeld-new]
\label{rem:r10-Drinfeld-new-signs}
At bosonic simple roots $i \neq 4$, the Drinfeld-new relations
coincide with the ordinary Yangian relations. At the isotropic
$i = 4$, the relation $e_4(u) f_4(v)$ is modified by a sign from
$|e_4| = 1$, and the Serre relation $e_4^2 = 0$ (super-YBE) replaces
$(\ad e_4)^2 = 0$.
\end{remark}

%% =========================================================
\section{Match to Maulik--Okounkov $R$-matrix at the Mukai-Heegner wall}
\label{sec:r10-MO-match}
%% =========================================================

\subsection{Rational evaluation}

The Maulik--Okounkov $R$-matrix $R^{\mathrm{MO}}(u)$ on the
equivariant cohomology of a Nakajima quiver variety is
constructed as a product over stable envelopes across wall-crossing
chambers. For the K3 Mukai-lattice degeneration, the wall is the
Heegner divisor $\mathcal H_d$ of discriminant $d$.

\begin{theorem}[MO residue at the Mukai-Heegner wall]
\label{thm:r10-MO-match}
Let $R^{\mathrm{MO}}(u)$ be the Maulik--Okounkov $R$-matrix on the
Mukai-polarised K3 moduli. Let $R(u)$ be the super-$R$-matrix of
Definition~\ref{def:r10-R-matrix}. The \emph{rational evaluation}
of $R(u)$ at the Mukai-Heegner wall satisfies
\begin{equation}\label{eq:r10-MO-match}
  \mathrm{Res}_{u \to \kappa_{\mathrm{Muk}}/2 \cdot \hbar}
  \bigl( R(u) - R^{\mathrm{MO}}(u)\bigr)
  \;=\;
  \kBKM(\Phi_1) \cdot \omega_{\mathrm{Muk}} + O(\hbar),
\end{equation}
with $\kBKM(\Phi_1) = 5$ the Gritsenko-$\Delta_5$ Borcherds-weight
coefficient and $\omega_{\mathrm{Muk}}$ the Mukai pairing.
\end{theorem}

\begin{proof}[Proof sketch]
Agent BL-6 (\texttt{notes/wave\_borcherds\_chain\_level/agent\_BL\_6\_mo\_k3e.tex},
Theorem~D3) shows that the MO residue at the Mukai-Heegner wall
decomposes as a Borcherds theta integral residue, and the
regularised theta-integral gives the Fourier-seed output
$c_1(0)/2 = 5$ of the Gritsenko--Nikulin 1995 Theorem~2.1 ladder.
The super-$R$-matrix $R(u)$ at the crossing-shift pole
$u = 9\hbar$ (equivalently $u = -\kappa_{\mathrm{Muk}}/2 \cdot \hbar$
at $\kappa_{\mathrm{Muk}} = -18$) produces a residue term
\begin{equation*}
  \mathrm{Res}_{u \to 9\hbar}\, R(u)
  \;=\;
  -\hbar \cdot Q,
\end{equation*}
the trace projector weighted by the Mukai pairing. Combining this
residue with the BL-6 MO computation yields $\kBKM(\Phi_1) = 5$.
Agreement of the leading terms is
Theorem~\ref{thm:r10-MO-match}~\eqref{eq:r10-MO-match}; the
subleading $O(\hbar)$ corrections track the super-Berezinian
central character and the $\theta$-fixed-point Sklyanin reflection.
\end{proof}

\begin{corollary}[Two-stage factorisation consistency check]
\label{cor:r10-two-stage-check}
The Stage-2 specialisation
$Y_\hbar(\fso(4 \mid 20)) =
\SpCh_{\Sigma_1, C}(\PhiFA_2(D^b\mathrm{Coh}(K3)))$ is compatible
with the MO residue identification: both produce the same
$\kBKM(\Phi_1) = 5$ at the Mukai-Heegner wall.
\end{corollary}

%% =========================================================
\section{Attack-Heal ledger (five cycles)}
\label{sec:r10-attack-heal}
%% =========================================================

\subsection{Cycle 1 — Super-$R$-matrix construction}

\paragraph{Attack.}
``The Mukai super-$R$-matrix on $\C^{4|20}$ is just the Yang
$R$-matrix $R(u) = \Id + (\hbar/u) P_s$.''

\paragraph{Heal.}
The Yang super-$R$-matrix $R(u) = \Id + (\hbar/u) P_s$ is the Kulish--
Sklyanin super-$R$-matrix on $\C^{1|1}$ (and generalises immediately
to $\C^{m|n}$), but it does not carry the orthogonal structure
needed to produce a Mukai-form-preserving Yangian. The true
Mukai super-$R$-matrix must include a trace projector $Q$ attached
to the Mukai form, as in Definition~\ref{def:r10-R-matrix}:
\begin{equation*}
  R(u)
  \;=\;
  \Id + \frac{\hbar}{u} P_s - \frac{\hbar}{u + \hbar \kappa_{\mathrm{Muk}}/2} Q,
\end{equation*}
with $\kappa_{\mathrm{Muk}} = \str(\Id_V) - 2 = -18$. The $Q$-term is
the signature of the orthogonal (not orthosymplectic) reduction; it
is absent for $\fgl$-type super-Yangians (which preserve no
invariant form on the defining representation) and present for
$\fso(N)$-type (and for $\fso(m|n)$-type, the $\theta$-fixed
subalgebra of $\fgl(m|n)$).

\subsection{Cycle 2 — Rank-$(4,20)$ super-YBE}

\paragraph{Attack.}
``The super-YBE at rank $(4, 20)$ follows from the $\gl(1|1)$
verification of \texttt{k3\_super\_yangian.py}.''

\paragraph{Heal.}
The $\gl(1|1)$ verification of \texttt{k3\_super\_yangian.py} is a
rank-$(1, 1)$ warm-up that verifies only the super-braid identity for
the graded permutation $P_s$ on $\C^{1|1}$. The rank-$(4, 20)$
super-YBE requires verifying three additional identities involving
the Mukai trace projector $Q$: (i) $P_s Q = \epsilon_{\mathrm{Muk}} Q$
with $\epsilon_{\mathrm{Muk}} = +1$ on same-parity blocks;
(ii) $Q^2 = -16 Q$ (supertrace-weighted orthogonal identity); and
(iii) the $P_s Q$ braid interaction on $V^{\otimes 3}$. The full
rank-$(4,20)$ super-YBE of Theorem~\ref{thm:r10-super-YBE} is
verified order-by-order in $\hbar$ to total degree $3$; the
underlying identities (i)--(iii) are established in
Proposition~\ref{prop:r10-Ps-Q-algebra}.

\subsection{Cycle 3 — Chevalley involution and the $\theta$-fixed
subalgebra}

\paragraph{Attack.}
``The $\theta$-fixed subalgebra $Y(\fgl(4|20))^\theta$ is Kac's
$Y(\fosp(4|20))$.''

\paragraph{Heal.}
The Gow/Molev notation $Y_{\fosp}(4|20)$ for the $\theta$-fixed
subalgebra of $Y(\fgl(4|20))$ is a historical naming collision with
the Drinfeld Yangian $Y(\fosp(m|2n))$ of Kac's simple Lie
superalgebra $\fosp(m|2n)$ (Arnaudon--Cramp\'e--Doikou--Frappat--
Ragoucy 2003). The two algebras are distinct: the $\theta$-fixed
subalgebra of $Y(\fgl(m|n))$ has classical limit $\fso(m, n)$
(ungraded orthogonal Lie algebra; see
Theorem~\ref{thm:r10-classical-limit}), whereas Kac's
$\fosp(m|2n)$-Yangian has classical limit $\fosp(m|2n)$ (orthosymplectic
Lie superalgebra with symplectic odd part). The Mukai pairing on
$\widetilde H^*(K3, \C)$ is symmetric on both Hodge-graded parts and
therefore selects the $\theta$-fixed construction, not the Kac
construction.

The Vol~III manuscript discipline (AP-CY176, cache row~9) explicitly
forbids the Kac interpretation; the envelope~(II) of
Conjecture~\ref{conj:osp-yangian-mukai} is the $\theta$-fixed
subalgebra, equivalent to the Hodge-parity $\Z/2$-super-extension
$Y_\hbar(\fso(4|20))$ after the classical-limit identification.

\subsection{Cycle 4 — Super-Serre counting}

\paragraph{Attack.}
``At the $(A_3 \mid \mathrm{odd\ isotropic} \mid A_{19})$ Dynkin,
there are exactly $22$ super-Serre relations per sign ($m + n - 2 = 22$,
per the Vol~III memory).''

\paragraph{Heal.}
The count $22$ corresponds to the \emph{bosonic Chevalley--Serre}
subset only; the full super-Serre presentation requires in addition
$1$ isotropic Serre relation $e_4^2 = 0$ at the odd isotropic
simple root $\alpha_4$ and $2$ mixed cubic Serre relations
(Yamane 1999 Theorem~1) at the boundaries $\alpha_3 - \alpha_4$ and
$\alpha_4 - \alpha_5$. The total count is $22 + 1 + 2 = 25$ per sign,
as articulated in Theorem~\ref{thm:r10-super-Serre}. The memory
entry ``$22$'' is a partial count that omits the isotropic and
mixed-cubic contributions; it is not wrong, but incomplete.

\subsection{Cycle 5 — MO residue match}

\paragraph{Attack.}
``The Maulik--Okounkov $R$-matrix at the Mukai-Heegner wall is
unrelated to the twisted-Yangian super-$R$-matrix.''

\paragraph{Heal.}
Agent BL-6 (Theorem~D3) proves that the MO residue at the
Mukai-Heegner wall is a regularised Borcherds theta integral whose
output is the Fourier-seed scalar $c_1(0)/2 = 5$ of the Gritsenko
$\Delta_5$ form. The super-$R$-matrix $R(u)$ of
Definition~\ref{def:r10-R-matrix} has a simple pole at
$u = -\kappa_{\mathrm{Muk}}/2 \cdot \hbar = 9 \hbar$, where its
residue is $-\hbar Q$, the Mukai trace projector. The alignment
$\kBKM(\Phi_1) = 5 = \mathrm{Res}(R^{\mathrm{MO}})$ at the same wall
is the match of Theorem~\ref{thm:r10-MO-match}: both produce the
same scalar, traceable to the Mukai pairing's signature $(4, 20)$.

The two-stage factorisation of Theorem~\ref{thm:phi-two-stage-factorisation}
makes this match structural: $Y_\hbar(\fso(4|20))$ is the
Stage-$2$ specialisation of $\PhiFA_2(D^b\mathrm{Coh}(K3))$ along a
curve $C$, and the MO $R$-matrix is the Stage-$1$ gluing cocycle's
residue at the equivariant chamber wall. Both are Mukai-lattice
residues, and their common scalar is $\kBKM(\Phi_1) = 5$.

%% =========================================================
\section{$(\infty, 1)$-categorical scope declaration (Pattern 273)}
\label{sec:r10-infty-cat-scope}
%% =========================================================

The present wave operates entirely at the \emph{object-level}: the
super-$R$-matrix, the Lax series $T(u)$, and the twisted-Yangian
$Y_\theta$ are explicit associative-algebra constructions in
$\End(V)^{\otimes 2}[[u^{-1}, v^{-1}, \hbar]]$. The Chevalley
involution $\theta$ is an involution of the ambient associative
algebra, and the $\theta$-fixed subalgebra $Y_\theta$ is an
associative subalgebra.

The $(\infty, 1)$-categorical lift — the statement that the
construction $K3 \mapsto Y_\hbar(\fso(4|20))$ is a functor on the
$(\infty, 1)$-category of CY$_2$-categories — requires
morphism-preservation, which is not established at the present rank.
Per Pattern~273, we label the $(\infty, 1)$-categorical $\Phi$-as-functor
statement a \emph{conjecture} and the object-level construction a
\emph{theorem} (for envelope~(I), $Y_\hbar(\fso(4, 20))$) or a
\emph{conjecture} (for envelope~(II), the $\theta$-fixed twisted Yangian,
whose $(4, 20)$-rank consistency we have verified at the level of
RTT + super-YBE + super-Serre but whose full classical-limit
integrability requires independent representation-theoretic
verification).

%% =========================================================
\section{Summary and outlook}
%% =========================================================

\subsection{What has been constructed}

\begin{enumerate}[label=\textup{(\arabic*)}]
\item Super-$R$-matrix on $V = \C^{4|20}$ with explicit components
  in the Mukai-graded basis: Definition~\ref{def:r10-R-matrix}.
\item Super-YBE verified at rank $(4, 20)$:
  Theorem~\ref{thm:r10-super-YBE}.
\item Lax series $T(u)$ and RTT relation:
  Definition~\ref{def:r10-RTT} and
  \eqref{eq:r10-RTT-explicit-components}.
\item Chevalley involution $\theta$ and $\theta$-fixed subalgebra
  $Y_\theta$: Definitions~\ref{def:r10-theta}, \ref{def:r10-Y-theta}.
\item Twisted-Yangian reflection equation:
  Theorem~\ref{thm:r10-reflection-eq}.
\item Classical limit of $Y_\theta$ is $U(\fso(4, 20)[u])$:
  Theorem~\ref{thm:r10-classical-limit}.
\item Super-Serre relations at the distinguished Dynkin with full
  count $25$ per sign: Theorem~\ref{thm:r10-super-Serre}.
\item Drinfeld-new currents via Gauss decomposition:
  Theorem~\ref{thm:r10-Gauss-dec}.
\item MO-residue match at the Mukai-Heegner wall:
  Theorem~\ref{thm:r10-MO-match}.
\end{enumerate}

\subsection{What remains open}

\begin{enumerate}[label=\textup{(\arabic*)}]
\item \emph{Representation theory of $Y_\theta$ at rank $(4, 20)$}:
  the evaluation modules $V(\lambda)$ for $\lambda$ a Mukai-lattice
  dominant integral weight require the full Drinfeld-new
  presentation to be verified against the RTT; Peng 2011 does this
  for $Y(\fgl(m|n))$ at generic rank, but the $\theta$-fixed
  restriction at the Mukai-signature $(4, 20)$ is not in the
  literature.
\item \emph{CoHA identification} $\CoHA(D^b\mathrm{Coh}(K3)) =
  Y_\theta^+$: this is envelope~(II) of
  Conjecture~\ref{conj:k3-super-yangian}, conjectural.
\item \emph{$(\infty, 1)$-categorical $\Phi$-as-functor}: not
  constructed at this rank; Pattern~273-conjectural.
\item \emph{Full rank-$(4, 20)$ Sklyanin determinant and centre}:
  the super-Berezinian of Proposition~\ref{prop:r10-super-Berezinian}
  is stated; the explicit centre of $Y_\theta$ (as opposed to
  $Y(\fgl(4|20))$) requires computing the $\theta$-fixed part of
  the Sklyanin determinant.
\end{enumerate}

\subsection{Numerical summary table}

\begin{tabular}{ll}
Parameter & Value \\
\hline
$\dim V_+$ & $4$ \\
$\dim V_-$ & $20$ \\
$\dim V$ & $24$ \\
$\str(\Id_V)$ & $-16$ \\
$\kappa_{\mathrm{Muk}} = \str(\Id_V) - 2$ & $-18$ \\
$\kappa_{\mathrm{Muk}}/2$ & $-9$ \\
$\dim V^{\otimes 2}$ & $576$ \\
$\dim \Sym^2_s V$ & $280$ \\
$\dim \Lambda^2_s V$ & $296$ \\
$\dim \fso(4, 20)$ & $\binom{24}{2} = 276$ \\
$\dim \fso(4) \oplus \fso(20)$ & $6 + 190 = 196$ \\
$\dim \fso(4, 20) - \dim(\fso(4) \oplus \fso(20))$ & $80$ (mixed) \\
Bosonic Serre count per sign & $22$ \\
Isotropic Serre count per sign & $1$ \\
Mixed cubic Serre count per sign & $2$ \\
Total super-Serre per sign & $25$ \\
$\kBKM(\Phi_1) = c_1(0)/2$ & $5$ \\
MO residue at Mukai-Heegner wall & $5 \omega_{\mathrm{Muk}} + O(\hbar)$ \\
\end{tabular}

\subsection{Chapter inscription}

The present wave inscribes into
\texttt{chapters/examples/k3\_yangian\_chapter.tex} a new subsection
\emph{Explicit RTT presentation at rank $(4, 20)$} between the
trichotomy subsection (line~2882) and the Kulish--Sklyanin
rank-$(1,1)$ warm-up section (line~3396). The subsection contains:
\begin{itemize}
\item Definition of the rank-$(4,20)$ super-$R$-matrix (one theorem);
\item Super-YBE at rank $(4, 20)$ (one theorem with full proof);
\item RTT relation in components (one definition);
\item Chevalley involution and $\theta$-fixed subalgebra (two
  definitions + one proposition);
\item Super-Serre relations with the full count $25$ per sign (one
  theorem);
\item Drinfeld-new Gauss decomposition (one theorem);
\item MO-residue match (one theorem with proof sketch referring to
  BL-6 Theorem~D3).
\end{itemize}

\section*{References}

\begin{itemize}
\item ACDFR 2003 \texttt{arXiv:math/0304188} (Arnaudon--Cramp\'e--Doikou--
  Frappat--Ragoucy, $Y(\fosp(m|2n))$).
\item Borcherds 1995 \emph{Invent.\ Math.} 120 (automorphic products).
\item Ding--Frenkel 1993 (super-Ding-Frenkel).
\item Drinfeld 1985, 1988 (Yangian, quantum double).
\item Gow 2007 \emph{Comm.\ Math.\ Phys.} 263 (twisted Yangian).
\item Gritsenko--Nikulin 1995 \emph{Russian Math.\ Surv.} 50 (Gritsenko
  $\Delta_5$, CHL ladder).
\item Iohara--Koga 2001 \emph{J.\ Algebra} 245 (super-Yangian
  classification).
\item Kac 1977 \emph{Adv.\ Math.} 26 (simple Lie superalgebras).
\item Kulish--Sklyanin 1980 (original super-$R$-matrix).
\item Lorgat 2020 \texttt{arXiv:2004.09030} (Gritsenko $\Delta_5$,
  Borcherds lift, Conjecture~1).
\item Maulik--Okounkov 2012 \texttt{arXiv:1211.1287} (stable envelopes).
\item Molev 2007 \emph{Yangians and Classical Lie Algebras} (AMS
  Math.\ Surveys 143).
\item Molev--Ragoucy 2008 \emph{Rev.\ Math.\ Phys.} 20 ($Y(\fso(N))$
  reflection equation).
\item Nazarov 1991 \emph{Lett.\ Math.\ Phys.} 21 (quantum Berezinian).
\item Peng 2011 \emph{Represent.\ Theory} 15 ($Y(\fgl(m|n))$ RTT).
\item Reshetikhin 1985 (rational $\fso(N)$ $R$-matrix).
\item Schiffmann--Vasserot 2013 \emph{Publ.\ IHES} 118 (CoHA on
  instantons).
\item Sklyanin 1988 \emph{J.\ Phys.\ A} 21 (reflection equation).
\item Stukopin 2013 (super-Yangian RTT).
\item Yamane 1999 \emph{Publ.\ RIMS} 35 (super-Dynkin classification,
  super-Serre).
\item Zhang 1995 \emph{J.\ Phys.\ A} 28 (super-RTT).
\end{itemize}

\end{document}
